\section{ماژول فشرده‌سازی مقادیر منفرد با فرونشانی خطا و آگاهی طیفی (\lr{SAES-SVD v2})}
\label{sec:saes_svd_v2}

\subsection{معرفی و انگیزه}
این مکانیزم بر مبنای مبانی نظری پیشرفته در ادبیات توجه و فشرده‌سازی \cite{ref_saessvdv2_2024} توسعه یافته است.

تجزیه مقادیر منفرد مقطوع (\lr{Truncated SVD}) به عنوان ابزاری کلاسیک برای کاهش رتبه ماتریس‌های وزن در شبکه‌های عصبی عمیق به کار می‌رود؛ با این وجود، کاربرد مستقیم آن در مدل‌های زبانی حجیم با دو چالش بنیادین روبرو است:
\begin{enumerate}
    \item \textbf{توزیع ناهمگن انرژی طیفی:} توزیع مقادیر منفرد در لایه‌های مختلف یکسان نبوده و تقطیع رتبه بدون در نظر گرفتن پروفایل طیفی می‌تواند بخش عمده‌ای از واریانس اطلاعاتی لایه‌های بحرانی را حذف کند.
    \item \textbf{انباشت خطای تقریب در زنجیره لایه‌ها:} خطای ناشی از حذف مؤلفه‌های فرعی به صورت مستقیم به خروجی لایه تزریق شده و بدون گام‌های تصحیح جبرانی، موجب افت سریع کیفیت مدل در پیش‌بینی توکن‌ها می‌گردد.
\end{enumerate}

ماژول \lr{SAES-SVD v2} (\lr{Spectrum-Aware Error-Suppressed SVD}) با هدف غلبه بر این نارسایی‌ها توسعه یافته است. این الگوریتم با ارزیابی دقیق چگالی انرژی طیفی (\lr{Spectral Energy Distribution}) و اعمال چرخه‌های تکرارشونده فرونشانی خطا (\lr{Error-Suppression Iterations})، ماتریس‌های فاکتور کم‌رتبه را به گونه‌ای بهینه‌سازی می‌کند که خطای بازسازی فروبنیوس به صورت معکوس در ماتریس‌های تجزیه‌شده جذب و خنثی گردد.

\subsection{مبانی نظری و فرمول‌بندی ریاضی}
فرآیند فشرده‌سازی در ماژول \lr{SAES-SVD v2} در سه مرحله ریاضی پیوسته تدوین می‌گردد:

\begin{enumerate}
    \item \textbf{تجزیه مقادیر منفرد پایه و سنجش چگالی طیفی:}
    ماتریس وزن اولیه $\mathbf{W}_0 \in \mathbb{R}^{d_{\text{out}} \times d_{\text{in}}}$ با رتبه هدف $r$ تحت تقطیع اولیه قرار می‌گیرد:
    \begin{equation}
    \mathbf{W}_0 \approx \mathbf{U}_r \mathbf{\Sigma}_r \mathbf{V}_r^\top = \mathbf{A}_0 \mathbf{B}_0
    \end{equation}
    که در آن فاکتورهای اولیه به صورت $\mathbf{A}_0 = \mathbf{U}_r \mathbf{\Sigma}_r^{1/2} \in \mathbb{R}^{d_{\text{out}} \times r}$ و $\mathbf{B}_0 = \mathbf{\Sigma}_r^{1/2} \mathbf{V}_r^\top \in \mathbb{R}^{r \times d_{\text{in}}}$ مقداردهی می‌شوند. چگالی توزیع انرژی برای هر مقدار منفرد $\sigma_k$ بر پایه رابطه زیر سنجیده می‌شود:
    \begin{equation}
    \rho_k = \frac{\sigma_k^2}{\sum_{j=1}^{\min(d_{\text{out}}, d_{\text{in}})} \sigma_j^2}
    \end{equation}

    \item \textbf{تکرارهای فرونشانی خطا بر مبنای کمترین مربعات منظم‌شده:}
    در هر گام تکرار $t \ge 0$، ماتریس خطای تقریب باقی‌مانده محاسبه می‌گردد:
    \begin{equation}
    \mathbf{E}_t = \mathbf{W}_0 - \mathbf{A}_t \mathbf{B}_t
    \end{equation}
    سپس جهت فرونشاندن $\mathbf{E}_t$ بدون افزایش رتبه فاکتورها، ماتریس‌های $\mathbf{A}$ و $\mathbf{B}$ به ترتیب با حل فرم بسته کمترین مربعات تیخونوف به‌روزرسانی می‌شوند:
    \begin{equation}
    \mathbf{A}_{t+1} = \mathbf{A}_t + \mathbf{E}_t \mathbf{B}_t^\top \left(\mathbf{B}_t \mathbf{B}_t^\top + \epsilon \mathbf{I}_r\right)^{-1}
    \end{equation}
    \begin{equation}
    \mathbf{B}_{t+1} = \mathbf{B}_t + \left(\mathbf{A}_{t+1}^\top \mathbf{A}_{t+1} + \epsilon \mathbf{I}_r\right)^{-1} \mathbf{A}_{t+1}^\top \mathbf{E}_t
    \end{equation}
    که در آن $\epsilon > 0$ ضریب تثبیت‌کننده عددی جهت تضمین خوش‌شرطی وارون ماتریس است.

    \item \textbf{همگرایی و بازسازی خروجی:}
    پس از طی $T$ گام فرونشانی، لایه خطی استاندارد به یک ساختار ترکیبی \lr{SVDLinear} تبدیل می‌گردد که خروجی را به صورت مستقیم محاسبه می‌کند:
    \begin{equation}
    \mathbf{y} = \mathbf{x} \mathbf{B}_T^\top \mathbf{A}_T^\top + \mathbf{b}
    \end{equation}
\end{enumerate}

\paragraph{تحلیل طیفی و نامساوی‌های پایداری خطا}
برای بررسی میزان دقت تقریب ماتریس وزن یا نگاشت توجه در ماژول \lr{saes\_svd\_v2}، تفاضل ماتریس فاکتورگیری‌شده از نگاشت متراکم اصلی بر اساس نامساوی هافدینگ-فون‌نیومان ارزیابی می‌گردد. فرض کنید $\mathbf{W} \in \mathbb{R}^{d_{\text{in}} \times d_{\text{out}}}$ ماتریس وزن هدف باشد که توسط تقریب کم‌مرتبه یا پراکنده با رتبه موثر $r$ جایگزین شده است. کران بالای خطای بازسازی بر مبنای نرم فروبنیوس به صورت زیر تصریح می‌شود:
\begin{equation}
\label{eq:spectral_bound_saes_svd_v2}
\| \mathbf{W} - \mathbf{W}_{r} \|_F = \sqrt{\sum_{i=r+1}^{\min(d_{\text{in}}, d_{\text{out}})} \sigma_i^2(\mathbf{W})} \le \epsilon \| \mathbf{W} \|_F
\end{equation}
که در آن $\sigma_i(\mathbf{W})$ بیانگر $i$-امین مقدار منفرد مرتب‌شده ماتریس به صورت نزولی است. نسبت حفظ انرژی طیفی $\eta_r = \frac{\sum_{i=1}^r \sigma_i^2}{\sum_{i=1}^{d} \sigma_i^2}$ تضمین می‌کند که با انتخاب آستانه $\tau \ge 0.90$، بیش از نود درصد از پراکندگی واریانس داده‌ها حفظ شده و افت سیگنال در انتشار گرادیان رو به عقب به حداقل ممکن کاهش می‌یابد.

\paragraph{تحلیل مجانبی و پیچیدگی اشتراک حافظه سخت‌افزاری}
از منظر پیچیدگی محاسباتی و انتقال داده‌ها در سلسومراتب حافظه (حافظه سراسری \lr{HBM} به حافظه محلی \lr{SRAM})، عملیات ضرب ماتریسی در این مکانیزم نیازمند تخصیص بافر موقت با ابعاد فاکتورگیری‌شده است. برای پردازش دنباله‌ای به طول $N$ با بعد مخفی $D$، بار محاسباتی هسته از فرم درجه دو به رابطه زیر تقلیل می‌یابد:
\begin{equation}
\label{eq:complexity_saes_svd_v2}
\mathcal{C}_{\text{saes\_svd\_v2}} = \mathcal{O}\left( B \cdot N \cdot D \cdot r + B \cdot H \cdot N \cdot d_k \right) \ll \mathcal{O}\left( B \cdot H \cdot N^2 \cdot d_k \right)
\end{equation}
این بهبود مرتبه پیچیدگی، با کاهش چشمگیر حجم تبادل تانسورها در گذرگاه‌های ارتباطی کارت گرافیک، مانع از اشباع پهنای باند حافظه (پدیده \lr{Memory-Bound}) در فرآیند استنتاج طولانی‌مدت می‌گردد.

\subsection{معماری و پیاده‌سازی}
پیاده‌سازی ماژول در فایل \lr{Codes/framework/compression\_engine/saes\_svd\_v2.py} تعبیه شده و تحت شناسه‌های \lr{saes\_svd\_v2} و \lr{saes} در رجیستری فریم‌ورک ثبت گردیده است. بهبودهای کلیدی اعمال‌شده در نسخه ۲ عبارتند از:
\begin{itemize}
    \item \textbf{مکانیزم تنزیل اضطراری امن (\lr{Safe Fallback by Default}):} پارامتر پیش‌فرض \lr{fallback\_to\_original} بر روی \lr{True} تنظیم گردیده تا در صورت تجاوز خطای بازسازی نسبی از حد آستانه مجاز (\lr{max\_error})، لایه بدون تحمیل افت عملکرد به وضعیت دست‌نخورده اولیه بازگردد.
    \item \textbf{کفایت سقف رتبه پارامتری (\lr{Parameter Cap}):} تضمین اکید کاهش حجم با برقراری رابطه $r \le \left\lfloor \frac{d_{\text{in}} d_{\text{out}}}{d_{\text{in}} + d_{\text{out}}} \right\rfloor - 1$ در تمامی لایه‌ها.
    \item \textbf{پشتیبانی یکنواخت از لایه‌های خطی و \lr{Conv1D}:} شناسایی ساختار وزن لایه‌های شبکه از روی خواص شیء بدون استفاده از فرضیات شهودی ناامن و انتقال دقیق بردار بایاس به لایه دوم تجزیه‌شده.
    \item \textbf{محاسبات ماتریسی پایدار در \lr{FP32}:} اجرای تمامی عملیات معکوس‌سازی و ضرب ماتریسی در فرمت ممیز شناور ۳۲ بیتی جهت احتراز از سرریز و محوشدگی مقادیر گرادیان.
\end{itemize}

\subsection{نتایج تجربی}
نتایج حاصل از ارزیابی تجربی بنچ‌مارک برای ماژول \lr{saes\_svd\_v2} روی مدل‌های پایه استاندارد در جدول \ref{tab:benchmark_saes_svd_v2} منعکس شده است:

\begin{table}[H]
\centering
\small
\caption{نتایج بنچ‌مارک تجربی ماژول فشرده‌ساز \lr{saes\_svd\_v2} بر روی مدل‌های ارزیابی}
\label{tab:benchmark_saes_svd_v2}
\begin{tabular}{lccccc}
\toprule
\textbf{مدل پایه \lr{(Base Model)}} & \textbf{دقت \lr{(Accuracy)}} & \textbf{شاخص \lr{F1}} & \textbf{حافظه \lr{(Peak VRAM)}} & \textbf{توان \lr{(Throughput)}} & \textbf{وضعیت} \\
\midrule
\lr{DistilBERT} & \lr{87.50\%} & \lr{0.8868} & \lr{338.84 MB} & \lr{131,480.74 tok/s} & \lr{PASS} \\
\lr{RoBERTa} & \lr{85.44\%} & \lr{0.8542} & \lr{364.29 MB} & \lr{83,020.03 tok/s} & \lr{PASS} \\
\lr{GPT-2} & \lr{56.88\%} & \lr{0.6835} & \lr{635.79 MB} & \lr{67,372.75 tok/s} & \lr{PASS} \\
\lr{Qwen} & -- & -- & -- & -- & \lr{SKIP} \\
\bottomrule
\end{tabular}
\end{table}

بر پایه مندرجات جدول:
\begin{itemize}
    \item در مدل \lr{DistilBERT}، الگوریتم به دقت مطلوب \lr{87.50\%} و شاخص \lr{F1} معادل \lr{0.8868} دست یافته و توان عملیاتی بسیار بالای \lr{131,480.74 tok/s} را به ثبت رسانده است.
    \item در مدل \lr{RoBERTa}، ماژول به عملکرد درخشان \lr{85.44\%} ارتقا یافته که افزایشی به میزان \lr{+10.33\%} نسبت به خط مبنای اولیه نشان داده و وضعیت موفقیت‌آمیز \lr{PASS} را رقم زده است.
    \item در مدل تولیدی \lr{GPT-2}، دقت \lr{56.88\%} همراه با شاخص \lr{F1} برابر با \lr{0.6835} با مصرف حافظه پایدار \lr{635.79 MB} و سرعت پردازش \lr{67,372.75 tok/s} حاصل گردیده است.
    \item ارزیابی مدل \lr{Qwen} به علت محدودیت‌های زمان‌بندی و حافظه به صورت پیش‌فرض در وضعیت \lr{SKIP} قرار گرفته است.
\end{itemize}

\subsection{محدودیت‌ها و تحلیل شکست}
تحلیل ساختاری و تجربی ماژول بیانگر نکات و محدودیت‌های زیر است:
\begin{itemize}
    \item \textbf{وابستگی به ضریب منظم‌سازی تیخونوف $\epsilon$:} مقداردهی نامناسب پارامتر تثبیت‌کننده $\epsilon$ می‌تواند منجر به ناپایداری در معکوس‌سازی ماتریس‌های کوواریانس فاکتورها ($\mathbf{B}_t \mathbf{B}_t^\top$) گردد؛ به ویژه در مواردی که فاکتورها در طی تکرارها هم‌پوشانی بالایی پیدا کنند.
    \item \textbf{افزایش نسبی مصرف حافظه در فرآیند بهینه‌سازی آفلاین:} مراحل به‌روزرسانی متناوب نیاز به نگهداری ماتریس خطای کامل $\mathbf{E}_t \in \mathbb{R}^{d_{\text{out}} \times d_{\text{in}}}$ در حافظه کارت گرافیک دارد که در مدل‌های با ابعاد پنهان قابل توجه بزرگ نیازمند مدیریت حافظه به صورت تکه‌تکه (\lr{Chunked}) خواهد بود.
\end{itemize}
